\documentclass[11pt,a4paper,openany]{book}
\input{t04_manual_style.tex}
\begin{document}
\frontmatter
\begin{titlepage}
\centering
\vspace*{1.1cm}
{\Huge\bfseries\color{navy}T04.P01\par}
\vspace{0.35cm}
{\LARGE\bfseries Manual operativo de teoría para problemas\par}
\vspace{0.2cm}
{\Large Divisibilidad, congruencias y técnicas híbridas\par}
\vspace{0.7cm}
{\large Semana 01 · Tema 04 · Oposiciones de Matemáticas\par}
\vfill
\begin{tcolorbox}[colback=lightblue,colframe=navy,width=.88\textwidth,arc=2mm,boxrule=.7pt]
\centering
Sistema de tres niveles para reconocer, resolver, justificar y transferir las técnicas necesarias en P2, P8, P10, P12, P15, P20, P24 y P25, con recursos audiovisuales selectivos y rutas de repaso.
\end{tcolorbox}
\vfill
{\large David Cortés Chaparro\par}
{\small Proyecto OPOSICIONES · versión 2.0 · agosto de 2026\par}
\end{titlepage}

\chapter*{Nota de alcance, fuentes y uso de los vídeos}
\addcontentsline{toc}{chapter}{Nota de alcance, fuentes y uso de los vídeos}
Este manual no sustituye al tema escrito T04.02 ni al tema esencial T04.03. Su función es distinta: desarrollar la teoría que permite \emph{resolver problemas nuevos}, elegir una técnica con criterio y justificar todos los pasos. Por ello incorpora resultados y procedimientos que serían excesivos en una redacción de tema, pero resultan decisivos en la prueba práctica.

El contenido parte de los materiales del proyecto disponibles en la sesión: T04.02 Tema de examen modular v4, T04.03 Tema esencial v5, T04.08 Selección semanal de 8 problemas, guía didáctica de aritmética, catálogo de las cuatro hojas de Ballesteros, relación de aritmética en problemas de oposición y mapa maestro T04 v2. Las ampliaciones operativas estándar se presentan como desarrollo matemático complementario, no como transcripción literal de una única fuente.

Los vídeos son \textbf{apoyos opcionales}. Se han seleccionado por su correspondencia temática y por la información pública disponible en sus páginas. El manual sigue siendo autosuficiente: ningún resultado depende de verlos. La rutina correcta es intentar primero, identificar el bloqueo, ver el recurso concreto, cerrar el vídeo y repetir el ejercicio sin ayuda.

\begin{nivelA}
El Nivel A se domina durante la primera semana: definiciones, hipótesis, procedimientos y ejemplos directamente conectados con los ocho problemas seleccionados.
\end{nivelA}
\begin{nivelB}
El Nivel B se trabaja en los repasos de 7 y 30 días: variantes, métodos alternativos y redacción completa.
\end{nivelB}
\begin{nivelC}
El Nivel C se activa durante el año: transferencia a problemas no vistos, generalizaciones y conexiones híbridas.
\end{nivelC}

\chapter*{Cómo utilizar el manual}
\addcontentsline{toc}{chapter}{Cómo utilizar el manual}
Cada capítulo contiene teoría mínima, señales de reconocimiento, procedimiento, ejemplos, errores frecuentes y conexión explícita con los problemas. La secuencia recomendada es:
\begin{enumerate}
\item leer la teoría activa y declarar las hipótesis;
\item cerrar el manual y reconstruir el procedimiento;
\item rehacer el ejemplo sin mirar;
\item resolver el problema vinculado;
\item anotar el primer paso que no surgió de forma natural;
\item repetir solo ese paso y volver a resolver el problema completo.
\end{enumerate}

\begin{ruta}[· Preguntas antes de calcular]
\begin{enumerate}
\item ¿La afirmación es de divisibilidad, resto, existencia, unicidad o clasificación?
\item ¿El módulo es primo, potencia de primo o producto de factores coprimos?
\item ¿La base es coprima con el módulo? Si no lo es, Euler o Fermat pueden no aplicarse.
\item ¿La expresión factoriza en enteros consecutivos, diferencias de cuadrados o potencias primas?
\item ¿Hay un número finito de estados o clases residuales que permita palomar o periodicidad?
\item ¿La congruencia obtenida es condición suficiente o solo necesaria?
\end{enumerate}
\end{ruta}

\tableofcontents
\mainmatter
\include{chapters/ch01}
\include{chapters/ch02}
\include{chapters/ch03}
\include{chapters/ch04}
\include{chapters/ch05}
\include{chapters/ch06}
\include{chapters/ch07}
\include{chapters/ch08}
\include{chapters/ch09}
\include{chapters/ch10}
\include{chapters/ch11}
\include{chapters/ch12}
\include{chapters/ch13}
\include{chapters/ch14}
\include{chapters/ch15}
\include{chapters/ch16}
\include{chapters/ch17}
\include{chapters/ch18}
\include{chapters/ch19}
\include{chapters/ch20}
\include{chapters/ch21}
\include{chapters/ch22}
\include{chapters/ch23}
\include{chapters/ch24}

\backmatter
\chapter*{Bibliografía y materiales del proyecto}
\addcontentsline{toc}{chapter}{Bibliografía y materiales del proyecto}
\begin{itemize}
\item Boyer, C. B. y Merzbach, U. C. \emph{Historia de la matemática}. Contexto funcional de Euclides y Gauss.
\item Hunacek, M. \emph{Introduction to Number Theory}. División euclídea, Bézout, diofánticas y congruencias.
\item Ivorra Castillo, C. \emph{Introducción a la teoría algebraica de números}. Inversos, congruencias lineales y TCR.
\item Apostol, T. M. \emph{Introducción a la teoría analítica de números}. Funciones aritméticas y multiplicatividad.
\item Ballesteros. Hojas H02 y H11 de aritmética y congruencias, según los materiales del proyecto.
\item Materiales del proyecto: T04.02, T04.03, T04.08, guía didáctica de aritmética, catálogo de cuatro hojas y mapa maestro T04 v2.
\end{itemize}

\begin{alerta}[· Control de versión]
Este manual debe actualizarse solo cuando un simulacro o un problema nuevo revele una carencia concreta. No se amplía por acumulación bibliográfica, sino por necesidad demostrada en la resolución.
\end{alerta}
\end{document}
